\chapter*{РЕАЛИЗАЦИЯ}
\addcontentsline{toc}{chapter}{РЕАЛИЗАЦИЯ}

На листинге 1 приведена реализация функции генерации ключей RSA.\\
\includelistingpretty
    {task-01.cpp}
    {}
    {Генерация ключей RSA}

Функция generate\_keys генерирует пару ключей для алгоритма RSA. Сначала генерируются два различных простых числа p и q в диапазоне 10000-50000. Затем вычисляется модуль $n = p \times q$ и функция Эйлера $\phi(n) = (p-1)(q-1)$. Публичная экспонента устанавливается равной 65537 (стандартное значение), если она взаимно проста с $\phi(n)$. Приватная экспонента d вычисляется как обратное к e по модулю $\phi(n)$ с использованием расширенного алгоритма Евклида.

На листинге 2 приведена реализация функции создания электронной подписи.\\
\includelistingpretty
    {task-02.cpp}
    {}
    {Создание электронной подписи}

Функция sign\_file создаёт электронную подпись для произвольного файла. Сначала файл полностью считывается в память. Затем вычисляется хэш-функция DJB2 от содержимого файла, получается 32-битное значение. Хэш приводится к диапазону модуля операцией остатка от деления. Подпись создаётся путём возведения хэша в степень d по модулю n (шифрование приватным ключом). Выходной файл формируется в формате: 4 байта размера подписи, 8 байт самой подписи, затем оригинальные данные файла.

На листинге 3 приведена реализация функции проверки электронной подписи.\\
\includelistingpretty
    {task-03.cpp}
    {}
    {Проверка электронной подписи}

Функция verify\_file проверяет подлинность электронной подписи. Подписанный файл считывается и проверяется корректность формата (минимальный размер, корректность размера подписи). Из файла извлекаются: размер подписи, сама подпись и оригинальные данные. Вычисляется хэш извлечённых данных. Подпись расшифровывается публичным ключом (возведение в степень e по модулю n), что даёт хэш-значение, которое было зашифровано при создании подписи. Если вычисленный хэш совпадает с расшифрованным, подпись считается верной и оригинальный файл сохраняется. В противном случае подпись признаётся недействительной.

На листинге 4 приведена реализация вспомогательных криптографических функций.\\
\includelistingpretty
    {task-04.cpp}
    {}
    {Вспомогательные функции RSA}

Функция mod\_pow выполняет быстрое модульное возведение в степень методом <<возведения в квадрат и умножения>>. Используется 128-битная арифметика (\_\_uint128\_t) для предотвращения переполнения при умножении 64-битных чисел. Функция compute\_hash реализует хэш-алгоритм DJB2 — простую, но эффективную хэш-функцию для демонстрационных целей. Функция gcd вычисляет наибольший общий делитель методом Евклида. Функция mod\_inverse вычисляет обратное по модулю с использованием расширенного алгоритма Евклида, необходимое для вычисления приватной экспоненты d.
